%% ════════════════════════════════════════════════════════════════════════════
\section{Literature Review and Related Work}
\label{sec:literature}
%% ════════════════════════════════════════════════════════════════════════════

\subsection{Smart Classroom Systems: State of the Art}
\label{subsec:lit_smart}

Smart classroom systems have evolved significantly over the past decade
from simple RFID card readers to multi-sensor IoT platforms.  Early
approaches such as~\cite{rfid_attendance_2018} demonstrated that RFID
tags could replace paper registers, reducing attendance time to under
two minutes per lecture.  However, RFID systems require students to
actively tap a reader, introducing a proxy-attendance vulnerability and
necessitating additional hardware distributed to each student.

Subsequent work explored QR-code-based attendance~\cite{qr_attend_2020},
which requires only a smartphone but shares the proxy problem and depends
on stable Internet connectivity for code rotation.  Camera-based
approaches eliminate student action entirely: the system identifies
occupants passively, closing the proxy loophole and reducing
professor workload to zero for the attendance step.

A recurring limitation across these systems is the absence of an
integrated environmental monitoring layer and of any AI-driven academic
risk analysis.  They treat attendance as the sole data point, ignoring
classroom conditions that correlate with student engagement and comfort.

\subsection{Face Recognition in Attendance Systems}
\label{subsec:lit_face}

Face recognition for attendance has been studied extensively using
conventional CNN-based pipelines~\cite{facenet_2015, deepface_2021}.
FaceNet~\cite{facenet_2015} introduced the 128-dimensional embedding
space with triplet-loss training, enabling sub-centimetre face-to-face
distance computation via cosine similarity.  DeepFace~\cite{deepface_2021}
unifies multiple model backends (Facenet, ArcFace, VGG-Face) behind a
single Python API, significantly reducing integration effort.

Our system uses DeepFace with the Facenet backend.  Enrollment
generates a 128-dimensional float32 embedding averaged from up to five
reference images per student; recognition applies cosine-distance
thresholding ($d < 0.40$) at 2~fps.  This threshold is tighter than
the DeepFace default ($0.60$) to reduce false-positive attendance
records on a shared-CPU Raspberry Pi~4B.

Prior work~\cite{rpi_face_2022} showed that Facenet inference on an
RPi~4B reaches 1--3~fps under normal lighting, consistent with our
empirical observations.  The primary tradeoff is latency: the 2~fps
ceiling means a 30-student class requires up to 15~seconds to scan
all faces in one frame pass, motivating the snapshot-per-cycle
attendance model described in Section~\ref{subsec:recognition}.

\subsection{IoT Protocols for Educational Environments}
\label{subsec:lit_mqtt}

Three lightweight messaging protocols are commonly evaluated for IoT
sensor networks: MQTT~\cite{mqtt_spec_5}, CoAP~\cite{rfc7252_coap},
and HTTP polling.  MQTT's publish-subscribe architecture decouples
producers (ESP32 sensor nodes) from consumers (backend services)
without requiring the sensor to know the consumer's address.  Its
binary framing and persistent TCP connection produce less per-message
overhead than HTTP, and its broker-mediated delivery survives transient
client disconnections through configurable quality-of-service (QoS)
levels.

QoS~0 (fire-and-forget) is used throughout this system.  For sensor
readings published every 5~seconds, an occasional dropped message
causes at most a 10-second gap in the time-series record---acceptable
for environmental monitoring.  Upgrading to QoS~1 would guarantee
delivery at the cost of increased broker state and acknowledgement
round-trips, which is unnecessary on a local network with sub-1~ms
latency.

\subsection{Local LLM Inference at the Edge}
\label{subsec:lit_llm}

Recent advances in model quantisation have made it feasible to run
billion-parameter language models on commodity CPU hardware.  The
phi3-mini model~\cite{phi3_mini_2024} (3.8~B parameters) from Microsoft
achieves strong reasoning quality with a 4k-token context window,
suitable for generating analytical prose from structured attendance data.
Ollama~\cite{ollama_2024} provides a local model server with a simple
REST API, eliminating the need for an API key or cloud subscription.

The rationale for a local LLM over a cloud API is threefold: (i)~student
attendance data should not leave the institution's network boundary;
(ii)~cloud API costs scale with usage and require internet access during
every pipeline run; and (iii)~a local model is always available,
decoupling the AI layer from internet connectivity.

The primary limitation is inference speed: on the RPi~4B CPU,
phi3-mini requires 10--30~seconds per student explanation.  The
at-risk pipeline batches all enrolled at-risk students into a single
pipeline run with one LLM call per student, and a Redis-based cooldown
(10-minute TTL) prevents redundant re-computation on each page
refresh~\cite{redis_docs}.

\subsection{LMS Integration}
\label{subsec:lit_moodle}

Moodle~\cite{moodle_docs} is the dominant open-source Learning Management
System in higher education.  It exposes a comprehensive REST API
authenticated via token, enabling programmatic creation and update of
grade and attendance records.  Automated Moodle sync has been
demonstrated in related work~\cite{moodle_auto_2021} as an effective way
to eliminate manual data entry and ensure that the official academic
record reflects real-time classroom events.

Our implementation calls the Moodle REST API on session end, writing
presence/absence records for each enrolled student.  A Redis retry
queue handles transient Moodle unavailability, with retry jobs running
on a 10-minute scheduler cycle.

\subsection{Summary and Positioning}
\label{subsec:lit_summary}

Table~\ref{tab:related_work} compares this system with four categories
of existing approaches across seven distinguishing features.  Our
system is the only one in this comparison that combines passive face
recognition with real-time environmental monitoring, local LMS
synchronisation, and on-premises AI-driven at-risk detection---all
without cloud dependency.

\begin{table*}[t]
\centering
\caption{Comparison of Classroom Management Approaches}
\label{tab:related_work}
\begin{tabular}{@{}lccccc@{}}
\toprule
\textbf{Feature}
  & \textbf{RFID~\cite{rfid_attend_2018}}
  & \textbf{QR-Code~\cite{qr_attend_2020}}
  & \textbf{CNN Face~\cite{rpi_face_2022}}
  & \textbf{Cloud IoT~\cite{cloud_iot_class_2023}}
  & \textbf{This Work} \\
\midrule
Attendance method        & RFID card  & QR scan   & Face rec.  & Face rec. & Face rec. \\
Proxy-attendance proof   & No         & No        & Yes        & Yes       & Yes       \\
Environment monitoring   & No         & No        & No         & Yes       & Yes       \\
LMS synchronisation      & No         & No        & No         & Partial   & Yes       \\
AI risk detection        & No         & No        & No         & No        & Yes       \\
Cloud-free deployment    & Yes        & No        & Yes        & No        & Yes       \\
Real-time dashboard      & No         & No        & No         & Yes       & Yes       \\
Open-source stack        & Partial    & Partial   & Partial    & No        & Yes       \\
\bottomrule
\end{tabular}
\end{table*}
